%-------------------------------------------------------------------------------
\section{Related Work}\label{sec:related}

\textbf{Low-Rank Transformer.}
Standard Transformer attention~\cite{OriginalTransformer} has $\mathcal{O}(n^2)$ complexity. Sparse patterns~\cite{SparseTransformer,Longformer} and memory compression~\cite{CompressiveTransformer} reduce cost but sacrifice full attention. Low-rank methods restore linear complexity while approximating the full attention map: Linformer~\cite{Linformer} projects keys and values into a lower-dimensional space; Performer~\cite{Performer} approximates softmax via random-feature kernelization; and Nystr\"omformer~\cite{Nystromformer} reconstructs the attention map from $m$ landmark tokens via the Nystr\"om method, reporting higher accuracy than Performer on standard benchmarks with $\mathcal{O}(n)$ complexity.

\textbf{zk-SNARKs for Matrix Multiplication.}
Thaler~\cite{Thaler} reduces matrix-multiplication verification to $\mathcal{O}(n^2)$ via Sumcheck; zkCNN~\cite{zkCNN} and zkML~\cite{zkML} adopt this for neural-network layers, the latter using Freivalds' algorithm~\cite{Freivalds}. vCNN~\cite{vCNN} represents convolutions as Quadratic Polynomial Programs; pvCNN~\cite{pvCNN} extends this to Quadratic Matrix Programs, casting a convolution layer as a single matrix product. zkLLM~\cite{zkLLM} achieves prover complexity linear in tensor size via specialized sumcheck protocols, building on SafetyNets~\cite{SafetyNets}. zkMatrix~\cite{zkMatrix} uses Bulletproofs-based acceleration for $\mathcal{O}(n^2)$ prover time with $\mathcal{O}(\log n)$ verification. However, all these approaches retain $\mathcal{O}(n^2)$ prover cost in the matrix dimension~\cite{Thaler,zkCNN,vCNN,pvCNN,zkML,zkMatrix} or, in the case of zkLLM~\cite{zkLLM}, remain linear in the size of the $n \times n$ attention tensor and hence quadratic in the sequence length, and nonlinear functions such as softmax require approximation~\cite{zkVC}. None of these lines of work exploits the low-rank structure of attention on the proving side; \name closes this gap by verifying a \nystrom-approximated attention map directly.
